\section*{Builder: construction plan}
\noindent Source: \texttt{structured/proposer.py}
\begin{lstlisting}
You are an expert at knowledge graph modeling with property graphs.
Design construction rules that turn the CSV files below into a graph serving the user's goal.

<goal>
{goal}
</goal>

The data profile below was computed exactly from the files. Trust it over your intuition:
- `is_unique` tells you which columns can be node keys. Never use a non-unique column as `unique_column`.
- `foreign_keys` lists columns whose values are contained in another file's unique column.
  Candidates with `name_match: false` on small integer columns are usually coincidences.{declared_rule}

<profile>
{profile}
</profile>

Modeling rules:
- A file with one unique identifier is a node. Its foreign key columns become relationships
  (source_file is that same file, from_column is its own key, to_column is the foreign key column).
- A file with no unique identifier and two foreign keys is a relationship file; its other columns
  are relationship properties.
- A file with no unique column whose rows describe things of their own, not only a link between two other
  files, is a node with `key_columns` (e.g. frames listed per hive, keyed by hive_id and frame_no): the
  fewest columns that together tell its rows apart and are never empty. Its `unique_column` is a new name
  for the key that code builds from them.
- Do not import foreign key columns as node properties.
- Set `name_column` to the column holding the name a person would write in a document
  (e.g. hive_name "Linden Hive", not a code like "hive_unit_07"); text is linked to nodes by it.
- Connect every pair of tables that the foreign keys join, directly or through a link table; tables that
  no key joins may stay apart. Skip files that are irrelevant to the goal.
- No two relationships between the same pair of labels may be inverses or synonyms of each other.

{feedback}
\end{lstlisting}

\section*{Builder: plan critic}
\noindent Source: \texttt{structured/proposer.py}
\begin{lstlisting}
You are reviewing a proposed knowledge graph construction plan.
The plan already passed mechanical checks (columns exist, keys are unique, whether one column or the
columns in key_columns, tables the keys join are connected),
so judge only the modeling:
- Could any node really be a relationship, or the reverse?
- Are relationship directions and types natural for the goal?
- Are obvious relationships from the foreign key candidates missing?
- Are any relationships redundant?
- Can the goal's typical questions be answered by traversing this schema?

Reply "retry" only for problems that would change the plan; otherwise "valid".

<goal>
{goal}
</goal>

<profile>
{profile}
</profile>

<plan>
{plan}
</plan>
\end{lstlisting}

\section*{Builder: text schema}
\noindent Source: \texttt{text/schema.py}
\begin{lstlisting}
You design the schema for extracting knowledge from unstructured text into a knowledge graph.

<goal>
{goal}
</goal>

The graph already contains this structured (domain) data. When the text talks about the same kind of
thing as a domain node, name the entity type like that node, so that code can link the two, and keep the
meaning of each type distinct from these:
<domain_graph>
{domain}
</domain_graph>

Representative text chunks:
<chunks>
{chunks}
</chunks>

Rules:
- Derive the entity types from the goal and the text: the things the text makes statements about that
  the goal cares about. Do not make types for ratings, dates or generic adjectives.
- Do not make an entity type for the source a statement comes from (a document, a report, a message, a
  post), and do not connect facts through such a source: where each fact comes from is recorded outside
  the schema. Connect the things a statement is about directly to each other.
- Propose fact types as (subject_type, PREDICATE, object_type) for every kind of claim the text makes
  about those things, not only the claims one question needs: a claim the schema cannot hold is lost.
- Describe aspects of the things, not only what goes wrong: what they are made of, how they are
  measured or rated, what they do well and what they do badly. Each extracted claim records whether it is
  positive, negative or neutral, so one fact type holds all three; do not split a fact type by tone.
- When the object of a claim is a number with a unit (a measure, a weight, a limit, a price), use the
  built-in object type {value_type}. It needs no entity type of its own.
- Give every entity type its identity. "keyed": each thing of the type is one record of the domain
  graph; name its labels in record_labels, and in key_attributes the record properties the text writes
  next to a name to tell two records of the same name apart. "individual": each thing is one particular
  person, organisation, place or object that the text names and the domain graph does not hold.
  "concept": the type names kinds, states or properties that many things share.
- Mark a fact type part_of when its claims state that the subject is one of the pieces the object is
  made of.
- Every fact type must reference only entity types you defined, and every entity type must be used by a
  fact type. Keep both lists small and precise (at most 12 entity types and 20 fact types).

{feedback}
\end{lstlisting}

\section*{Builder: claim extraction}
\noindent Source: \texttt{text/extraction.py}
\begin{lstlisting}
Extract facts from the text chunk as subject-predicate-object triples.

Allowed entity types:
{entity_types}

Allowed fact types (subject_type -[PREDICATE]-> object_type):
{fact_types}

Rules:
- Use only the entity types and fact types above. Skip anything that does not fit.
- Be exhaustive: one triple per distinct claim. A sentence that lists several claims gives one triple
  per item of the list. A mild or hedged claim ("a bit", "somewhat", "seems to", "more than expected")
  still counts. A claim about the thing the document is about as a whole goes on that thing. When one
  statement supports two of the allowed fact types, state both facts.
- `subject` and `object` are names exactly as written in the text, with only the words that name the thing:
  leave out a clause that says when or under which condition the claim holds ("when ...", "if ...",
  "after ..."), words of frequency ("sometimes", "often"), and the words that deny the fact or say it only
  may hold; they stay in the evidence and in the fields below. Never use pronouns.
- The chunk comes from the document named in <document>. When the text refers to the thing the document
  is about with a pronoun or a generic word ("it", "this kettle"), use the proper name from the document
  name as the entity name.
- `evidence` must be ONE contiguous quote copied verbatim from the chunk that supports the fact.
- `polarity` is the claim's own tone toward its subject: "positive" when it speaks well of it, "negative"
  when it reports a fault, a harm or dissatisfaction, "neutral" when it only states what is so.
- `time`: when the quote says when the claim holds or after how long ("since ...", "within ..."),
  copy those words verbatim from the quote; otherwise leave it empty. A condition is not a time.
- `truth`: "negated" when the text says the fact does not hold or did not happen, in whatever words:
  "the kettle never leaked" denies a leak, "the filter prevents scale" denies scale; otherwise "affirmed".
  Name the fact itself ("leak", not "never leaked"). Keep the denial in a name only when the name is itself
  the denied state; the fact is still negated and still has its `negation`: "the kettle will not switch
  off" can give the object "will not switch off" with truth "negated" and negation "will not".
- `negation`: for a negated fact, the words of the quote that deny it, copied verbatim ("never",
  "prevents"); otherwise leave it empty.
- `modality`: "possible" when the text says the fact may hold or happen, not that it does; "conditional"
  when it holds under a condition the text states, in whatever words ("if ...", "unless ...", "when ..."
  meaning whenever, "had the lid been shut"); otherwise "actual", also for what a thing is able to do
  ("holds two litres"). A fact with both a condition and a possibility is conditional.
- `hedge`: the words of the quote that say the fact only may hold, copied verbatim ("may", "could", "a
  risk of"); required for a possible fact, allowed for a conditional one, otherwise empty.
- `condition`: for a conditional fact, the words of the condition copied verbatim from the quote ("when
  the water boils"); otherwise leave it empty.
- When a fact type's object type is {value_type}, `object` is only the number and its unit, copied from
  the quote.
- Extract only what the text states. Do not infer. Return an empty list when nothing qualifies.

<document>{context}</document>
<chunk id="{chunk_id}">
{text}
</chunk>
\end{lstlisting}

\section*{Builder: mention pass}
\noindent Source: \texttt{text/mention\_pass.py}
\begin{lstlisting}
List the things this text names or talks about that are not listed yet: the entries an
index of the text would have, so that a reader looking for one of them finds this text.

In:
- named things: people, places, organisations, events, works, awards and identifiers;
- kinds the text says something about: an object or a piece of one; a state, condition or fault, a property
  included ("the brightness fades": brightness); an incident or event; a work done to a thing that makes,
  fits, repairs, cleans, tests, replaces or withdraws it (an installation, a recalibration, an inspection);
  a person or organisation known only by its role ("technician", "ferry operator");
- a fault, state or incident of a thing, or a work done to it, written as a verb, as the text writes it
  ("the shutter sticks": sticks; "the hull was repainted": repainted);
- each of these also when the text says it is absent or did not happen ("No condensation was found on the
  mirror": condensation, mirror).
Out:
- a describing word alone; a light or reporting verb such as "is", "has", "showed" or "explained"; a clause
  or evaluation that names no thing; a quantity, date, time or season; a pronoun;
- words that would fit any text, even as the subject or object of a sentence ("thing", "issue",
  "experience", "number", "time", "people");
- everyday acts of people, the ordinary use of a thing included, that are not a fault, an incident or a work
  done to the thing (contacting, filing, buying, reading, boarding a ferry, looking through a telescope);
- the text or document itself; its writer or reader referred to only as such ("the writer", "the person
  reporting", "the undersigned"); a form of address ("Madam", "Dear neighbours"), though a name in it counts;
- a title, role or common noun written next to a name or in apposition with it when it names that very
  thing: the name is the only entry, also when the name is already listed ("dome technician Edit Varga" and
  "Edit Varga, the dome technician": Edit Varga; "ferry T-4471": T-4471). A noun for the many things of a
  named model or class names other things and stays ("the Skylark 30 ferries": Skylark 30, ferries).

Copy each name verbatim from the text, as whole words and without "a", "an" or "the": the shortest span that
names the thing, keeping a describing word only when it tells the thing apart ("shutter motor", not "old
shutter motor in the dome"). List each thing once; when the text gives it two names ("Night Vision Unit
(NVU)"), list the first.

Give every thing two separate answers:
- mention_class, how the text refers to it: "particular" when the text calls one individual thing by its own
  name, a proper name or an identifier ("North Dome", "Lakeside Ferries", "Varga Prize", "T-4471"), also a
  unit of a named organisation called with that name ("Lakeside Ferries ticket office"); "kind" for
  everything else, also for a common noun that means one thing here ("the dome" is a kind even where it is
  North Dome, and so are "the ticket office" and "the morning crossing");
- type, what it is: the type below whose description fits it, whatever its mention_class; "{none}" when no
  type fits.

Types:
{types}

Already listed (do not list these again): {known}

Text ({document}):
{text}
\end{lstlisting}

\section*{Builder: record chooser}
\noindent Source: \texttt{resolution/record\_choice.py}
\begin{lstlisting}
A text names a thing "{name}". Is it the very same thing as one of the records below?

Where "{name}" is named ([document] sentence):
{lines}

Records (id, then what the data holds about it):
{records}

Choose a record if "{name}" is that record itself, perhaps written in other words: a plural, a fuller or
shorter name, the name with a describing word, a version of it (a variant of the same line, in another size
or form), or the name with a word for the side of it the sentence talks about, such as its design or how it
is made ("the brass focuser" and "the focuser design" are the record "Focuser"; "the Corvid ED" is the record
"Corvid").
Answer none if "{name}" is another thing that only shares words with a record: a piece of it, something made
for it or fixed to it, or a larger thing it belongs to ("the thread of the focuser" and "the focuser cap" are
not the record "Focuser").
Answer none if a word added to the record's name makes the name of another line sold under a name of its own
("the Corvid Voyager" is not the record "Corvid" when the Voyager is a line of its own).
Answer none if "{name}" means other things of the record's kind than the ones it stands for, such as earlier
ones the sentence sets apart from it ("the Corvids I owned before" are not the record "Corvid" whose data
gives this year's edition).
If "{name}" fits a whole and also pieces of that whole (the relations show which record is a piece of which),
it is the whole, unless the sentences name the piece. Answer none when several records still fit equally well.
If you choose a record, copy its id into record and copy one sentence from the lines above into quote,
verbatim. Otherwise answer none and leave quote empty.
\end{lstlisting}

\section*{Builder: individual identity}
\noindent Source: \texttt{resolution/individuals.py}
\begin{lstlisting}
Are A and B, both of type {etype}, the same particular person, organisation, place or
object? Two different ones may share a name, and one may be written in several ways.

Answer "same" only when the lines show that A and B are one: the same role, position, place, organisation,
date or event is stated for both. A similar name is not enough, and neither is naming the same record or
place: two different people can belong to one organisation. Answer "different" when the lines show two
different ones (a role, a place or a time that cannot hold for both). Answer "unsure" when the lines do not
settle it.
For example, "E. Varga (dome technician)" and "Edit Varga, who services the dome" can be "same"; two
"E. Varga" whose lines only both mention the observatory are "unsure".
If you answer "same", copy one of A's lines into quote_a and one of B's lines into quote_b, verbatim: the
lines that show it. Otherwise leave both quotes empty.

A: {a}
{record_a}Lines that name A, with their neighbours ([document] sentence; "*" marks a line that names A):
{lines_a}

B: {b}
{record_b}Lines that name B, with their neighbours ([document] sentence; "*" marks a line that names B):
{lines_b}

Records the texts of both A and B also name:
{shared}
\end{lstlisting}

\section*{Builder: concept identity}
\noindent Source: \texttt{resolution/concepts.py}
\begin{lstlisting}
Do these two names, both of type {etype}, name the same kind of thing?
{etype} names kinds, states or properties that many things share, so the names may come from different
documents: answer whether A and B are the same kind, even when they are mentioned in different documents.
Different sizes, models, materials, pieces or people are NOT the same. Answer conservatively.

A: {a}
B: {b}

Where A is mentioned ([document] sentence):
{ctx_a}

Where B is mentioned ([document] sentence):
{ctx_b}
\end{lstlisting}

\section*{Query engine: agent}
\noindent Source: \texttt{agent/prompt.py}
\begin{lstlisting}
You answer a question about one dataset. Tools search its knowledge graph, the text it was read from and its source tables; you call them one at a time, then answer from what they returned.

Rules:
- Each turn, either call one tool (`action` "call", `tool` its name, `arguments` a JSON object of its
  arguments) or answer (`action` "answer").
- Start from the nodes the question names (find_nodes) and walk from them; read the schema first when the
  question asks how many, which is largest or which records meet a condition.
- Take a count or a sum of entries from a tool's `measure`, with their conditions in `where`, never by
  counting listed rows yourself; to leave out the entries another result returned, add one `!=` condition
  per id it returned.
- Use only what the tools returned.
- Set `outcome` to what the results let you give:
  - "answer": they state the answer, and you cite them; a result showing there are none is an answer
    (`entities` []);
  - "insufficient_evidence": the dataset holds what the question is about, but no result states the answer;
  - "out_of_scope": the schema and your searches show the dataset does not hold what the question is about;
  - "clarify": the question can be read in several ways that would get different answers; list them in
    `readings` and ask in `text` which one is meant.
  With any outcome but "answer", leave `entities` and `number` empty and say why in `text`.
- Put the answer in the field its form calls for:
  - one number (a count, an amount, a measurement, a year; how many, how much, how far, which year): give
    it in `number`, also when it is stated as a bound or an estimate;
  - one or several named things, or values that are neither one number nor a date (who, which, where,
    what is the name or the value of ...; several years; a month alone or a time of day): list them in
    `entities`, each once, by its name or value alone, as the tools or the question give it;
  - anything else (what was said or done, why, whether, a date: a day of a month, with or without its
    year): answer in a short `text`.
  `text` may add to `entities` or `number`; it never replaces them.
- Write an entity by its name, never by a ref.
- Cite every id your answer relies on, exactly as a tool returned it: a passage by its chunk id with ONE
  contiguous quote copied verbatim from its text; a node, a claim or a query result by its id alone.
- Never repeat a call: its result is already below.

<question>{question}</question>

<tools>
{tools}
</tools>

<turns>
{turns}
</turns>

Tool calls left: {left}.
\end{lstlisting}

\section*{Query engine: closing prompt at a cap}
\noindent Source: \texttt{agent/prompt.py}
\begin{lstlisting}
You answer a question about one dataset from the tool results below. No tool can be called any more: reply now (`action` "answer") from these results alone.
- Set `outcome` to what the results let you give:
  - "answer": they state the answer, and you cite them; a result showing there are none is an answer
    (`entities` []);
  - "insufficient_evidence": the dataset holds what the question is about, but no result states the answer;
  - "out_of_scope": the schema and your searches show the dataset does not hold what the question is about;
  - "clarify": the question can be read in several ways that would get different answers; list them in
    `readings` and ask in `text` which one is meant.
  With any outcome but "answer", leave `entities` and `number` empty and say why in `text`.
- Put the answer in the field its form calls for:
  - one number (a count, an amount, a measurement, a year; how many, how much, how far, which year): give
    it in `number`, also when it is stated as a bound or an estimate;
  - one or several named things, or values that are neither one number nor a date (who, which, where,
    what is the name or the value of ...; several years; a month alone or a time of day): list them in
    `entities`, each once, by its name or value alone, as the tools or the question give it;
  - anything else (what was said or done, why, whether, a date: a day of a month, with or without its
    year): answer in a short `text`.
  `text` may add to `entities` or `number`; it never replaces them.
- Cite every id your answer relies on exactly as a tool returned it; a passage by its chunk id with ONE contiguous quote copied verbatim from its text.

<question>{question}</question>

<turns>
{turns}
</turns>
\end{lstlisting}

\section*{Query engine: passage reranker}
\noindent Source: \texttt{hybrid/passage\_rerank.py}
\begin{lstlisting}
You order text passages for a reader who must answer a question from them. Each passage is shown whole, under its document's name.

The question: {question}

The passages:
{passages}

Rules:
1. Order the passages by how directly each one states something the question asks for, in any wording: a passage that says it in other words counts as much as one that repeats the question's words.
2. A question can have several answers in several passages. Rank each of them high, not only the single best one.
3. Judge from the passages only.
4. In order, best first, give the id of every passage, such as "P3", each exactly once.
\end{lstlisting}
